\section{The GreenFaaS Algorithm}
\label{sec:algorithm}

This section presents the \GreenFaaS{} scheduler. \S\ref{sec:algorithm:design}
gives a high-level narrative. \S\ref{sec:algorithm:pseudocode} presents
the algorithm in pseudocode. \S\ref{sec:algorithm:sla} describes the
SLA-tiered policy. \S\ref{sec:algorithm:complexity} analyses time and
space complexity. \S\ref{sec:algorithm:estimation} discusses
arrival-rate estimation and jitter.

\subsection{Design Overview}
\label{sec:algorithm:design}

\GreenFaaS{} schedules each invocation independently at its arrival,
producing a decision $(r(i), s(i))$. The design has three layers:
\textbf{(1)} an arrival-rate estimator tracks per-(function,
home-region) arrival rates online with an EWMA; \textbf{(2)} a
\emph{region gate} applies Lemma~\ref{lemma:tradeoff} to each
candidate region given the current $\hat\lambda$ and present
carbon intensities, dropping regions that fail the gate not because
they could not save carbon \emph{now} but because routing this
function class to them on a sustained basis would lose carbon to
warm-pool idle energy; \textbf{(3)} a \emph{per-slot scorer}
enumerates a small grid of (region, start-time) pairs over the
admitted regions and the SLA-bounded deferral window, selecting the
pair with minimum expected carbon, with an idle-energy term that
charges deferred invocations for the additional warm-container
holding cost (\S\ref{sec:tradeoff:idle}).

These three layers are stateless across invocations except for the
arrival-rate estimator. There is no global queue, no
inter-invocation coordination, no pile-up risk from synchronised
defer slots (we use a per-function deterministic jitter to ensure
this; \S\ref{sec:algorithm:estimation}). The scheduler can run in
parallel across invocations with no locks beyond the rate update.

\subsection{Algorithm}
\label{sec:algorithm:pseudocode}

We present \GreenFaaS{} as Algorithm~\ref{alg:greenfaas} with
sub-procedures \textsc{LemmaGate} and \textsc{ScoreSlots}. We write
$\mathcal{S}(t)$ for the platform state at $t$: warm-pool counts
$W_{r,f}(t)$ and in-flight counts $N_r(t)$.

\begin{algorithm}[t]
\caption{\GreenFaaS$(i, \mathcal{S}, \hat{C})$}
\label{alg:greenfaas}
\begin{algorithmic}[1]
\State $\hat\lambda \gets \textsc{UpdateRate}(f, h, a(i))$
\If{$\ell(f) = \Interactive$} \Return $(h, a(i))$
\EndIf
\State $(\rho^*, \Delta^*) \gets \textsc{ClassLimits}(\ell(f))$
\State $\mathcal{R}_0 \gets \{r : \rho(h,r) \le \rho^* \wedge N_r/K_r < u_{\max}\}$
\State $\mathcal{R}^* \gets \{r \in \mathcal{R}_0 : \textsc{LemmaGate}(f, h, r, \hat\lambda, \ldots)\}$
\If{$\mathcal{R}^* = \emptyset$} \State $\mathcal{R}^* \gets \{h\}$
\EndIf
\State \Return $\textsc{ScoreSlots}(f, h, \mathcal{R}^*, \Delta^*, \hat{C}, \mathcal{S}, a(i))$
\end{algorithmic}
\end{algorithm}

\begin{algorithm}[t]
\caption{\textsc{LemmaGate}$(f, h, r, \hat\lambda, \mathcal{S}, \hat{C}, t)$}
\label{alg:lemmagate}
\begin{algorithmic}[1]
\If{$r = h$} \Return \textbf{true}
\EndIf
\State $C_r \gets \hat C_r(t),\;\; C_h \gets \hat C_h(t)$
\If{$C_r \ge C_h$} \Return \textbf{true} \Comment{admit; scorer compares instantaneous carbon}
\EndIf
\State $\alpha \gets \tau_c(f)/\tau_e(f)$
\State $\beta \gets P_i(f) T_w / (P_a(f)\tau_e(f))$
\State $\mathrm{ratio} \gets C_h / C_r$
\State $\beta_{\mathrm{crit}} \gets (\mathrm{ratio}-1)(1+\alpha)$
\If{$\beta \le \beta_{\mathrm{crit}}$} \Return \textbf{true} \Comment{Regime 1}
\EndIf
\State $\lambda^* \gets \textsc{SolveBreakEven}(\mathrm{ratio}, \alpha, \beta, T_w)$
\State \Return $\hat\lambda > \lambda^*$
\end{algorithmic}
\end{algorithm}

\begin{algorithm}[t]
\caption{\textsc{ScoreSlots}$(f, h, \mathcal{R}^*, \Delta^*, \hat{C}, \mathcal{S}, t_0)$}
\label{alg:scoreslots}
\begin{algorithmic}[1]
\State $(r^*, s^*, S^*) \gets (\bot, \bot, +\infty)$
\For{each $r \in \mathcal{R}^*$}
  \State $\delta \gets$ forecast step; $N \gets \lceil \Delta^*/\delta \rceil + 1$
  \State $\bar{C}_r \gets$ long-run mean of $C_r$
  \State $w_0 \gets \mathbf{1}[W_{r,f}(t_0) > 0]$
  \For{$k = 0, 1, \ldots, N-1$}
    \State $t_k \gets t_0 + k\delta + \text{jitter}(f, i, k)\cdot\delta$
    \State $\text{cold} \gets \neg(w_0 \wedge k=0)$
    \State $\text{eta} \gets (t_k - t_0) + \rho(h,r) + \text{cold}\cdot\tau_c(f) + \tau_e(f)$
    \If{$\text{eta} > D(f)$} \textbf{continue}
    \EndIf
    \State $C_k \gets \hat{C}_r(t_k)$
    \State $S \gets E_{\mathrm{exec}}(f,r)C_k + \text{cold}\cdot E_{\mathrm{cs}}(f,r)C_k + E_{\mathrm{idle}}(f,r,t_k-t_0)\bar{C}_r$
    \If{$S < S^*$} $(r^*, s^*, S^*) \gets (r, t_k, S)$
    \EndIf
  \EndFor
\EndFor
\If{$r^* = \bot$} \Return $(h, t_0)$
\EndIf
\State \Return $(r^*, s^*)$
\end{algorithmic}
\end{algorithm}

The three energy quantities are
$E_{\mathrm{exec}}(f, r) = P_a(f)\tau_e(f)\,\mathrm{PUE}_r / (3.6 \times 10^6)$,
$E_{\mathrm{cs}}(f, r) = P_a(f)\tau_c(f)\,\mathrm{PUE}_r / (3.6 \times 10^6)$, and
$E_{\mathrm{idle}}(f, r, \Delta t) = P_i(f)\Delta t\,\mathrm{PUE}_r / (3.6 \times 10^6)$,
all in kWh; multiplying by the appropriate carbon intensity gives
gCO$_2$eq. The use of $\bar{C}_r$ for the idle term (rather than
$\hat{C}_r(t_k)$) is the \S\ref{sec:tradeoff:idle} correction: idle
energy accrues \emph{over} the deferral window, not at a single
instant.

\paragraph{Operational semantics of the deferral idle term.}
GreenFaaS does \emph{not} reserve a dedicated warm container during a
deferral: a deferred invocation ($k > 0$) is conservatively treated
as a cold start (line~8, $\text{cold} \gets \neg(w_0 \wedge k{=}0)$),
because the platform makes no guarantee that a warm container will
survive the wait under contention. The
$E_{\mathrm{idle}}(f, r, t_k - t_0)\,\bar{C}_r$ term is therefore not
a literal per-invocation energy charge for a reserved container;
it is a \emph{conservative opportunity-cost penalty} that
approximates the additional warm-pool exposure the deferral induces
at the pool level (a delayed invocation lengthens the interval over
which the pool's warm capacity is held without serving work). This
conservatism is deliberate: it biases the scorer against deferral
unless the carbon saving clearly outweighs the penalty, which is
precisely the mechanism that yields the do-no-harm property
(\S\ref{sec:evaluation:fine-donoharm}). An alternative model that
reserves a warm container during deferral --- charging literal idle
energy and setting $\text{cold} \gets \neg(w_0 \wedge k\delta < T_w)$
--- would relax this conservatism; we adopt the no-reservation model
because it matches the default behaviour of FaaS platforms that do
not pin containers to pending requests.

\textsc{SolveBreakEven} solves~(\ref{eq:dichotomy}) for $u^* > 0$
by bisection on the strictly decreasing LHS, returning
$\lambda^* = u^*/T_w$. Convergence to $10^{-9}$ requires $\le 200$
iterations in practice.

\subsection{SLA-Tiered Policy}
\label{sec:algorithm:sla}

The class limits $(\rho^*, \Delta^*)$ in line 3 of
Algorithm~\ref{alg:greenfaas} encode SLA tiering. Defaults:

\begin{center}\small
\begin{tabular}{lrr}
\toprule
Class & $\rho^*$ (RTT) & $\Delta^*$ (defer) \\
\midrule
\Interactive  & 0\,ms      & 0\,s \\
\Deferrable   & 80\,ms     & $\min(60\,\mathrm{s}, D - \tau_e)$ \\
\Background   & $\infty$   & $D - \tau_e$ \\
\bottomrule
\end{tabular}
\end{center}

\Deferrable{} invocations are capped at 60\,s of deferral even when
their deadline is longer, which keeps the temporal-shift horizon
below the warm-pool TTL ($T_w = 600$\,s by default) and ensures the
idle-energy term remains a small correction rather than the dominant
cost. \Background{} invocations use the full deadline.

\subsection{Complexity}
\label{sec:algorithm:complexity}

Let $R = |\mathcal{R}|$, $N$ the number of forecast slots per region,
$B$ the bisection iterations in \textsc{SolveBreakEven}.
\textsc{UpdateRate} is $O(1)$. The candidate filtering is $O(R)$.
The region gate calls \textsc{LemmaGate} once per candidate, each
$O(B)$. \textsc{ScoreSlots} enumerates $O(R \cdot N)$ slots, each
$O(1)$ given precomputed $\bar{C}_r$. \textbf{Per-invocation time:}
$O(R \cdot (B + N))$. For default settings ($R \le 9$,
$N \le 12$ for \Deferrable, $N \approx 720$ for \Background), this
is well below 100\,$\mu$s per invocation in practice. \textbf{Per-invocation
space:} $O(R)$. The arrival-rate estimator uses
$O(|\mathcal{F}| \cdot R)$ state amortised across time.

\subsection{Arrival-Rate Estimation and Jitter}
\label{sec:algorithm:estimation}

\textsc{UpdateRate} maintains a per-(function, region) EWMA:
$\hat\lambda \leftarrow \gamma \cdot (t - t_{\mathrm{prev}})^{-1}
+ (1 - \gamma) \hat\lambda_{\mathrm{prev}}$ with smoothing
coefficient $\gamma = 0.2$ (we use $\gamma$ here to avoid collision
with $\alpha = \tau_c/\tau_e$ from \S\ref{sec:tradeoff}). Before
the first update, $\hat\lambda$ is initialised to a small
default (0.01/s) making the gate conservative for cold-start
functions.

Without jitter, all invocations of the same function in the same
deferral window would target the same defer slot --- exactly the
Wait-Awhile failure mode of synchronised pile-up. We apply a
deterministic per-function offset
$\mathrm{jitter}(f, i, k) = (H(f \oplus i) \bmod M) / M \cdot \phi$
for $k \ge 1$, where $H$ is a hash, $M = 65535$, and $\phi = 0.25$
is the max fraction of a slot width by which jitter shifts the
candidate time. The jitter is zero at $k=0$ so that immediate
execution remains a candidate.

\paragraph{Forecasts.} $\hat{C}$ may be the true carbon trace
(\textsf{perfect}), the true trace within a cutoff with Gaussian-noise
decay beyond (\textsf{1h}, \textsf{24h}), or a constant equal to each
region's historical mean (\textsf{none}). Forecast accuracy is a
sensitivity axis in \S\ref{sec:evaluation:forecast}; we will see that
\GreenFaaS{} is essentially forecast-invariant, an unintended but
welcome consequence of the design.
